\section{Overall Algorithm}
\label{sec:overall}

The message updates in Section~\ref{sec:algorithm}, with parameters
learned as described in Section~\ref{sec:learning}, define the
$L$-layer receiver. We now summarize its execution and analyze
the computational complexity.

\subsection{Message Scheduling}

From the received pilots, the receiver estimates the noise variance
by \eqref{eq:noise-mdl} and the phase increments by the fitting
procedure in Appendix~\ref{app:frontend}. It then constructs
$\widehat{\mathbf A}$ using \eqref{eq:effective-pilot}. The observations
and estimated noise variance are divided by $s$ and $s^2$,
respectively, as specified in Section~\ref{sec:system}. The resulting
$\mathbf A_{\rm c}$ and $\sigma_{\rm c}^2$ remain fixed throughout
the unfolded layers.

The first layer starts from the prior mean and variance
\begin{equation}
 \mathbf M_1=\mathbf0,\qquad w_1=w_{\rm ref}=\rho P_H.
 \label{eq:twc-initialization}
\end{equation}
Each layer updates the conditional channels and activity
probabilities, combines them into the effective channel estimate,
and returns an extrinsic message to the next layer. Let
$\widehat H_{k,L}$ denote the complex channel tensor obtained by
reshaping $\mu_{kL}$. After $L$ layers, the physical-scale outputs are
\begin{equation}
 \widehat H_k^{\rm out}=s\widehat H_{k,L},\qquad
 \widehat X_k^{\rm out}=s p_{kL}\widehat H_{k,L},\qquad
 \widehat\alpha_k=\mathbf1\{p_{kL}\geq1/2\}.
 \label{eq:twc-final-output}
\end{equation}
The activity probabilities $p_{kL}$ are also retained as soft outputs.
Table~\ref{tab:twc-algorithm} summarizes the complete receiver.

\begin{table}[t]
\centering
\caption{Unfolded receiver for joint activity detection and channel estimation}
\label{tab:twc-algorithm}
\begin{tabular}{p{0.93\linewidth}}
\toprule
\textbf{Input:} received pilots $\mathbf Y$, original pilot tensor $P$,
activity marginal $\rho$, trained layer parameters, and training
scales $s,P_H$.\\
\textbf{Preprocessing:} estimate $\widehat\sigma^2$ and
$\widehat{\bm\beta}$ by Appendix~\ref{app:frontend}; construct
$\mathbf A_{\rm c}=\widehat{\mathbf A}$.
Normalize $\mathbf Y$ and $\widehat\sigma^2$ by $s$ and $s^2$,
and set $\sigma_{\rm c}^2$ to the normalized noise estimate.\\
\textbf{Initialize:} use \eqref{eq:twc-initialization} and compute
the thin SVD of $\mathbf A_{\rm c}$.\\
\textbf{For} $\ell=1,\ldots,L$:\\
1. Compute the linear update and extrinsic query using
\eqref{eq:twc-linear-correction}--\eqref{eq:twc-query-calibrated}.\\
2. Estimate the channels and activity probabilities, and combine
them using \eqref{eq:twc-conditional-mean}, \eqref{eq:twc-activity},
and \eqref{eq:twc-estimate}.\\
3. If $\ell<L$, form the feedback message and damp it using
\eqref{eq:twc-residual-scale}--\eqref{eq:twc-damping}.\\
\textbf{Return:} the channel estimates and decisions in
\eqref{eq:twc-final-output}, together with $p_{kL}$.\\
\bottomrule
\end{tabular}
\end{table}

\subsection{Computational Complexity Analysis}

The online cost consists of phase and noise estimation, followed by
the unfolded receiver. For the phase estimator in
Appendix~\ref{app:frontend}, let $G_{\rm c}$ be the number of coarse
delay--Doppler path matrices, $B_\phi\geq Q$ the number of phase bins,
and $S_\phi$ the number of coordinate sweeps. Precomputing the path
matrices and their pilot responses costs
$O(G_{\rm c}N^3+KG_{\rm c}QN^2)$.
Each user search requires correlations at cost $O(G_{\rm c}QNM)$
and phase-bin evaluation by FFTs at cost
$O(G_{\rm c}MB_\phi\log B_\phi)$.

Let $E_\phi$ count path-matrix evaluations during alias ranking and
local refinement, including derivative evaluations up to a constant
factor. Each evaluation, pilot response, and gain fit costs
$O(N^3+QN^2+QNM)$. For $K\leq QN$, a joint gain refit after each
sweep costs $O(QNK^2+QNKM)$. Including the initial user search,
the phase-estimation cost is
\begin{equation}
 \begin{aligned}
 C_{\rm phase}=O\big(&G_{\rm c}N^3+KG_{\rm c}QN^2\\
 &+(S_\phi+1)KG_{\rm c}M(QN+B_\phi\log B_\phi)\\
 &+E_\phi(N^3+QN^2+QNM)+S_\phi(QNK^2+QNKM)\big).
 \end{aligned}
 \label{eq:twc-phase-complexity}
\end{equation}
Forming the spatial covariance and computing its eigenvalues for
noise estimation cost $O(QNM^2+M^3)$.

Within the unfolded receiver, the dominant computations are linear
estimation and network evaluation.
Since $\mathbf A_{\rm c}$ is fixed, one thin singular
value decomposition suffices to evaluate both the mean correction
and the trace in \eqref{eq:twc-linear-correction} across all layers.
With $r=\min(Q,J)$, this decomposition costs $O(QJr)$.
Each layer requires $O(QJC)$ operations for the residuals and
$O((Q+J)rC)$ for applying the singular vectors. If $F_s$ and $F_a$
denote the costs of one user's score-network and activity-classifier
evaluations, respectively, the cost per layer is
\begin{equation}
 C_{\rm layer}=O\left(QJC+(Q+J)rC+K(F_s+F_a)\right).
 \label{eq:twc-layer-complexity}
\end{equation}
The complete online cost is therefore
\begin{equation}
 O\left(C_{\rm phase}+QNM^2+M^3+QJr+L C_{\rm layer}\right).
 \label{eq:twc-complexity}
\end{equation}
For $Q\leq J$, the SVD and unfolded layers together cost
$O(Q^2J+L[QJC+K(F_s+F_a)])$, in addition to the phase and noise
estimators.
If the score network and activity classifier contain $P_s$ and $P_a$
parameters, the $L$ layers store $L(P_s+P_a+3)$ parameter scalars.
During inference, layers are evaluated sequentially. Excluding network
parameters and activations, the operator, SVD, and messages require
$O(QJ+(Q+J)r+JC+QC)$ storage. The phase estimator additionally
caches $O(G_{\rm c}N^2+KG_{\rm c}QN)$ entries, with
$O(G_{\rm c}B_\phi M)$ temporary FFT storage.
